\section{Native estimators}\label{sec:estimators_python_examples}
Below some basic examples on how to use certain estimators
are presented. For more comprehensive ones please visit
the \texttt{examples} subdirectory of each \gretlpy distribution.

In order to illustrate the \gretlpy estimating facilities, we:
\begin{enumerate}
	\item take it as given that \pythcode{formula: str} can
	take one of the following forms:
	\begin{itemize}
		\item a simple \proglang{R}-style formula that
		maps the model equation according to the following rule:
		\begin{center}
			\refstepcounter{equation}
			\hspace*{15pt}$\displaystyle y = \beta_0 + \sum_{j=1}^{k}\beta_jx_j + \varepsilon\;\rightarrow\;$%
			\pythcode{'y ~ const + x1 + x2 + ... + xk'}%
			\hfill(\theequation)\label{eq:Rstyle}
		\end{center}

		\item \gretl's default syntax, as described in \cite{gretl_manual},
		referred to here as the ``maps to'' syntax:
		\begin{center}
			\refstepcounter{equation}
			\hspace*{15pt}$\displaystyle y = \beta_0 + \sum_{j=1}^{k}\beta_jx_j + \varepsilon\;\rightarrow\;$%
			\pythcode{'y const x1 x2 ... xk'}%
			\hfill(\theequation)\label{eq:Gstyle}
		\end{center}
	\end{itemize}

	\item assume that \pythcode{varlist: list} is constructed
	according to the following pattern:
	\begin{center}
		\refstepcounter{equation}
		\hspace*{15pt}$\displaystyle y = \beta_0 + \sum_{j=1}^{k}\beta_jx_j + \varepsilon\;\rightarrow\;$%
		\pythcode{['y', 'const', 'x1', 'x2', ..., 'xk']}%
		\hfill(\theequation)\label{eq:GstyleList}
	\end{center}
	which is later transformed into \pythcode{'y const x1 x2 ... xk'}.
	Furthermore, \pythcode{[[var1, var2], [var3, var4]]} is always
	transformed into \pythcode{'var1 var2 ; var3 var4'}.
	Note that the \pythcode{list} in question may consist of either
	variable names or their IDs; see \citep[for details]{gretl_manual}.
\end{enumerate}
The complete list of available native estimators is as follows.%
\footnote{By ``native'' we mean an estimator implemented directly
	in \lgretl, as opposed to econometric methods (estimators,
	tests, etc.) available as contributed packages
	(see Section~\ref{sec:classpypackage}).}\newline
Note that all of these are factory functions, \textit{i.e.}, they
return an allocated \classpymodel instance.

\gretlfpycmd{ar()}%
\index{get\_data()}
\estdesc{To perform the generalized Cochrane-Orcutt iterative
	procedure, the following minimal code snippet can be used.
	The corresponding results are presented in Figure~\ref{fig:ar_results}.}
\begin{lstlisting}[style=pythonstylefoot]
import gretl

d = gretl.get_data('longley.csv')
m = gretl.ar('1 2 ; employ const prdefl gnp armfrc').fit()
print(m)
\end{lstlisting}
\generaldesc{See \texttt{`examples/estimators/ar.py'} for details.}

\examplefig{ar}{\label{fig:ar_results}}

\gretlfpycmd{ar1()}%
\index{get\_data()}
\estdesc{To perform feasible GLS estimates for a model in
	which the error term is assumed to follow a first-order
	autoregressive process, the following minimal code snippet
	can be used. The corresponding results are presented
	in Figure~\ref{fig:ar1_results}.}
\begin{lstlisting}[style=pythonstylefoot]
import gretl
  
d = gretl.get_data('longley.csv')
m = gretl.ar1('employ ~ const + prdefl + gnp').fit()
print(m)
\end{lstlisting}
\generaldesc{See \texttt{`examples/estimators/ar1.py'} for details.}

\examplefig{ar1}{\label{fig:ar1_results}}

\gretlfpycmd{arima()}%
\index{get\_data()}
\estdesc{To perform $ARIMA(1,1,1)$ estimation, the following
	minimal code snippet can be used:}
\begin{lstlisting}[style=pythonstylefoot]
import gretl

d = gretl.get_data('bjg.gdt')
m = gretl.arima(depvar='lg', order=(1,1,1)).fit()
print(m)
\end{lstlisting}
\generaldesc{See \texttt{`examples/estimators/arima.py'} for details.}

\examplefig{arima}{\label{fig:arima_results}}

\gretlfpycmd{biprobit()}%
\index{get\_data()}
\estdesc{To perform bivariate probit model estimation,
	the following minimal code snippet can be used. The
	corresponding results are presented in Figure~\ref{fig:biprobit_results}.}
\begin{lstlisting}[style=pythonstylefoot]
import gretl

d = gretl.get_data('greene25_1.gdt')
m = gretl.biprobit(
	[['anydrg', 'cardhldr', 'const', 'age', 'avgexp'],
	['const', 'age', 'income', 'ownrent', 'selfempl']]
).fit()
print(m)
\end{lstlisting}
\generaldesc{See \texttt{`examples/estimators/biprobit.py'} for details.}

\examplefig{biprobit}{\label{fig:biprobit_results}}

\gretlfpycmd{dpanel()}%
\index{get\_data()}
\estdesc{Carries out estimation of dynamic panel data models
	(that is, panel models including one or more lags of the
	dependent variable) using either the GMM-DIF or GMM-SYS method.
	The corresponding results are presented in Figure~\ref{fig:dpanel_results}.}
\begin{lstlisting}[style=pythonstylefoot]
import gretl

d = gretl.get_data('abdata.gdt')
m = gretl.dpanel(p=2,
	formula='n w k ; GMM(k,2,3) k GMM(w,2,3)').fit()
print(m)
\end{lstlisting}
\generaldesc{See \texttt{`examples/estimators/dpanel.py'} for details.}

\examplefig{dpanel}{\label{fig:dpanel_results}}

\gretlfpycmd{duration()}%
\index{get\_data()}
\estdesc{To perform duration model estimation, the following
	minimal code snippet can be used. The corresponding
	results are presented in Figure~\ref{fig:duration_results}.}
\begin{lstlisting}[style=pythonstylefoot]
import gretl

d = gretl.get_data('kennan.gdt')
m = gretl.duration('Ti ~ const + Prod').fit()
print(m)
\end{lstlisting}
\generaldesc{See \texttt{`examples/estimators/duration.py'} for details.}

\examplefig{duration}{\label{fig:duration_results}}

\gretlfpycmd{garch()}%
\index{get\_data()}
\estdesc{To perform $GARCH(1,1)$ model estimation, the
	following minimal code snippet can be used. The corresponding
	results are presented in Figure~\ref{fig:garch_results}.}
\begin{lstlisting}[style=pythonstylefoot]
import gretl

d = gretl.get_data('b-g.gdt')
m = gretl.garch('1 1 ; Y').fit()
print(m)
\end{lstlisting}
\generaldesc{See \texttt{`examples/estimators/garch.py'} for details.}

\examplefig{garch}{\label{fig:garch_results}}

\gretlfpycmd{heckit()}%
\index{get\_data()}
\estdesc{To perform Heckman-type selection model estimation,
	the following minimal code snippet can be used. The corresponding
	results are presented in Figure~\ref{fig:heckit_results}.}
\begin{lstlisting}[style=pythonstylefoot]
import gretl

d = gretl.get_data('mroz87.gdt')
m = gretl.heckit('WW const AX WE CIT ; LFP const WA FAMINC WE').fit()
print(m)
\end{lstlisting}
\generaldesc{See \texttt{`examples/estimators/heckit.py'} for details.}

\examplefig{heckit}{\label{fig:heckit_results}}

\gretlfpycmd{hsk()}%
\index{get\_data()}
\estdesc{This is applicable where heteroskedasticity is present
	in the form of an unknown function of the regressors which
	can be approximated by a quadratic relationship. The corresponding
	results are presented in Figure~\ref{fig:hsk_results}.}
\begin{lstlisting}[style=pythonstylefoot]
import gretl

d = gretl.get_data('data4-1')
m = gretl.hsk('price const sqft bedrms').fit()
print(m)
\end{lstlisting}
\generaldesc{See \texttt{`examples/estimators/hsk.py'} for details.}

\examplefig{hsk}{\label{fig:hsk_results}}

\gretlfpycmd{intreg()}%
\index{get\_data()}
\estdesc{Estimates an interval regression model. The corresponding
	results are presented in Figure~\ref{fig:intreg_results}.}
\begin{lstlisting}[style=pythonstylefoot]
import gretl

d = gretl.get_data('wtp')
d.sample('depvar == 1', restrict=True)
d.new_series('hi = bidl')
d.sample('depvar == 2', restrict=True, replace=True)
d.new_series('lo = bidl')
d.new_series('hi = bid1')
d.sample('depvar == 3', restrict=True, replace=True)
d.new_series('lo = bid1')
d.new_series('hi = bidh')
d.sample('depvar == 4', restrict=True, replace=True)
d.new_series(name='lo', value='bidh')
d.sample(full=True)

m = gretl.intreg("lo hi const age female income").fit()
print(m)
\end{lstlisting}
\generaldesc{See \texttt{`examples/estimators/intreg.py'} for details.}

\examplefig{intreg}{\label{fig:intreg_results}}

\gretlfpycmd{lad()}%
\index{get\_data()}
\estdesc{To calculate a regression that minimizes the sum
	of the absolute deviations of the observed from the fitted
	values of the dependent variable, the following minimal code
	snippet can be used. The corresponding results are
	presented in Figure~\ref{fig:lad_results}.}
\begin{lstlisting}[style=pythonstylefoot]
import gretl

d = gretl.get_data('longley.csv')
m = gretl.lad('employ ~ const + prdefl + gnp').fit()
print(m)
\end{lstlisting}
\generaldesc{See \texttt{`examples/estimators/lad.py'} for details.}

\examplefig{lad}{\label{fig:lad_results}}

\gretlfpycmd{logistic()}%
\index{get\_data()}
\estdesc{To perform Logistic regression estimation, the
	following minimal code snippet can be used. The corresponding
	results are presented in Figure~\ref{fig:logistic_results}.}
\begin{lstlisting}[style=pythonstylefoot]
import gretl

d = gretl.get_data('longley.csv')
m = gretl.logistic('employ ~ const + prdefl + gnp').fit()
print(m)
\end{lstlisting}
\generaldesc{See \texttt{`examples/estimators/logistic.py'} for details.}

\examplefig{logistic}{\label{fig:logistic_results}}

\gretlfpycmd{logit()}%
\index{get\_data()}
\estdesc{To perform logit regression estimation, the following
	minimal code snippet can be used. The corresponding
	results are presented in Figure~\ref{fig:logit_results}.}
\begin{lstlisting}[style=pythonstylefoot]
import gretl

d = gretl.get_data('greene22_2.gdt')
m = gretl.logit('Y ~ const + Z1 + Z2').fit()
print(m)
\end{lstlisting}
\generaldesc{See \texttt{`examples/estimators/logit.py'} for details.}

\examplefig{logit}{\label{fig:logit_results}}

\gretlfpycmd{midasreg()}%
\index{get\_data()}
\estdesc{Carries out least-squares estimation (either NLS or OLS,
	depending on the specification) of a MIDAS (Mixed Data Sampling)
	model. The corresponding results are presented in Figure~\ref{fig:midasreg_results}.}
\begin{lstlisting}[style=pythonstylefoot]
import gretl

d = gretl.get_data('gdp_midas.gdt')
d.new_list(name='dX', value='ld_payems*')
d.new_list('dXL = hflags(3, 11, dX)')
gretl.genr('theta = {1,5}')

m = gretl.midasreg(formula='ld_qgdp 0 ld_qgdp(-1) ; mdsl(dXL, 2, theta)').fit()
print(m)
\end{lstlisting}
\generaldesc{See \texttt{`examples/estimators/midasreg.py'} for details.}

\examplefig{midasreg}{\label{fig:midasreg_results}}

\gretlfpycmd{mpols()}%
\index{get\_data()}
\estdesc{To perform multiple precision floating-point
	arithmetic OLS estimation, the following minimal code
	snippet can be used. The corresponding results are presented
	in Figure~\ref{fig:mpols_results}.}
\begin{lstlisting}[style=pythonstylefoot]
import gretl

d = gretl.get_data('longley.csv')
m = gretl.mpols('employ ~ const + prdefl + gnp').fit()
print(m)
\end{lstlisting}
\generaldesc{See \texttt{`examples/estimators/mpols.py'} for details.}

\examplefig{mpols}{\label{fig:mpols_results}}

\gretlfpycmd{negbin()}%
\index{get\_data()}
\estdesc{To perform Negative Binomial model estimation,
	the following minimal code snippet can be used. The
	corresponding results are presented in Figure~\ref{fig:negbin_results}.}
\begin{lstlisting}[style=pythonstylefoot]
import gretl

d = gretl.get_data('rac3d.gdt')
m = gretl.negbin('DVISITS ~ const + SEX + AGE').fit()
print(m)
\end{lstlisting}
\generaldesc{See \texttt{`examples/estimators/negbin.py'} for details.}

\examplefig{negbin}{\label{fig:negbin_results}}

\gretlfpycmd{ols()}%
\index{get\_data()}
\estdesc{To perform OLS estimation, the following minimal
	code snippet can be used. The corresponding results are
	presented in Figure~\ref{fig:ols_results}.}
\begin{lstlisting}[style=pythonstylefoot]
import gretl

d = gretl.get_data(gretl.gretldir + '/data/longley.csv')
m = gretl.ols('employ ~ const + prdefl + gnp').fit()
print(m)
\end{lstlisting}
\generaldesc{See \texttt{`examples/estimators/ols.py'} for details.}

\examplefig{ols}{\label{fig:ols_results}}

\gretlfpycmd{panel()}%
\index{get\_data()}
\estdesc{To perform panel model estimation, the following
	minimal code snippet can be used. The corresponding
	results are presented in Figure~\ref{fig:panel_results}.}
\begin{lstlisting}[style=pythonstylefoot]
import gretl

d = gretl.get_data(gretl.gretldir + '/data/penngrow.gdt')
m = gretl.panel('Y ~ const + Y(-1) + X').fit()
print(m)
\end{lstlisting}
\generaldesc{See \texttt{`examples/estimators/panel.py'} for details.}

\examplefig{panel}{\label{fig:panel_results}}

\gretlfpycmd{poisson()}%
\index{get\_data()}
\estdesc{To perform Poisson regression estimation, the
	following minimal code snippet can be used. The corresponding
	results are presented in Figure~\ref{fig:poisson_results}.}
\begin{lstlisting}[style=pythonstylefoot]
import gretl

d = gretl.get_data(gretl.gretldir + '/data/rac3d.gdt')
m = gretl.poisson('DVISITS ~ const + SEX + AGE').fit()
print(m)
\end{lstlisting}
\generaldesc{See \texttt{`examples/estimators/poisson.py'} for details.}

\examplefig{poisson}{\label{fig:poisson_results}}

\gretlfpycmd{probit()}%
\index{get\_data()}
\estdesc{To perform probit regression estimation, the following
	minimal code snippet can be used. The corresponding
	results are presented in Figure~\ref{fig:probit_results}.}
\begin{lstlisting}[style=pythonstylefoot]
import gretl

d = gretl.get_data(gretl.gretldir + '/data/greene22_2.gdt')
m = gretl.probit('Y ~ const + Z1 + Z2').fit()
print(m)
\end{lstlisting}
\generaldesc{See \texttt{`examples/estimators/probit.py'} for details.}

\examplefig{probit}{\label{fig:probit_results}}

\gretlfpycmd{tobit()}%
\index{get\_data()}
\estdesc{To perform Instrumental Variables (IV) estimation
	using two-stage least squares (TSLS), the following
	minimal code snippet can be used. The corresponding
	results are presented in Figure~\ref{fig:tobit_results}.}
\begin{lstlisting}[style=pythonstylefoot]
import gretl

d = gretl.get_data('tobit.gdt')
m = gretl.tobit('
	apt ~ const + read + math + proggeneral + progvocational',
	llimit=400, rlimit=600.7).fit()
print(m)
\end{lstlisting}
\generaldesc{See \texttt{`examples/estimators/tobit.py'} for details.}

\examplefig{tobit}{\label{fig:tobit_results}}

\gretlfpycmd{tsls()}%
\index{get\_data()}
\estdesc{To perform Instrumental Variables (IV) estimation
	using two-stage least squares (TSLS), the following
	minimal code snippet can be used. The corresponding
	results are presented in Figure~\ref{fig:tsls_results}.}
\begin{lstlisting}[style=pythonstylefoot]
import gretl

d = gretl.get_data(gretl.gretldir + '/data/penngrow.gdt')
m = gretl.tsls('Y Y(-1) X ; 0 X Y(-2)').fit()
print(m)
\end{lstlisting}
\generaldesc{See \texttt{`examples/estimators/tsls.py'} for details.}

\examplefig{tsls}{\label{fig:tsls_results}}

\gretlfpycmd{var()}%
\index{get\_data()}
\estdesc{Sets up and estimates (using OLS) a vector autoregression
	(VAR). The corresponding results are presented in Figure~\ref{fig:var_results}.}
\begin{lstlisting}[style=pythonstylefoot]
import gretl

d = gretl.get_data('denmark.gdt')
m = gretl.var(order=1, ylist='LRM LRY').fit()
print(m)
\end{lstlisting}
\generaldesc{See \texttt{`examples/estimators/var.py'} for details.}

\examplefig{var}{\label{fig:var_results}}

\gretlfpycmd{vecm()}%
\index{get\_data()}
\estdesc{Sets up and estimates vector error-correction model (VECM).
	The corresponding results are presented in Figure~\ref{fig:vecm_results}.}
\begin{lstlisting}[style=pythonstylefoot]
import gretl

d = gretl.get_data('denmark.gdt')
m = gretl.vecm(order=1, rank=1, ylist='LRM LRY').fit()
print(m)
\end{lstlisting}
\generaldesc{See \texttt{`examples/estimators/vecm.py'} for details.}

\examplefig{vecm}{\label{fig:vecm_results}}

\gretlfpycmd{wls()}%
\index{get\_data()}
\estdesc{To perform weighted least squares (WLS) estimation,
	the following minimal code snippet can be used. The corresponding
	results are presented in Figure~\ref{fig:wls_results}.}
\begin{lstlisting}[style=pythonstylefoot]
import gretl

d = gretl.get_data(gretl.gretldir + '/data/longley.csv')
m = gretl.wls('year employ const prdefl gnp').fit()
print(m)
\end{lstlisting}
\generaldesc{See \texttt{`examples/estimators/wls.py'} for details.}

\examplefig{wls}{\label{fig:wls_results}}
